\documentclass[11pt,a4paper]{article}

% ---------------------------------------------------------
% PREAMBLE (stable, warning-free, Overleaf-safe)
% ---------------------------------------------------------
\usepackage[utf8]{inputenc}
\usepackage[T1]{fontenc}
\usepackage{lmodern}
\usepackage{amsmath, amssymb, amsfonts}
\usepackage{bm}
\usepackage{geometry}
\usepackage{graphicx}
\usepackage{hyperref}
\usepackage{cite}
\usepackage{authblk}
\usepackage{physics}
\usepackage{mathtools}

\geometry{margin=2.4cm}

\title{\textbf{Companion XC: End-to-End RDCC Inference Pipeline — From Simulations to Posteriors}}
\author{Michael Lehmann \\ \texttt{mi.lehmann@gmx.de}}
\date{April 2026}

\begin{document}
\maketitle

\begin{abstract}
The Relaxation-Driven Cyclic Cosmology (RDCC) introduces the relaxation parameter 
$\alpha$, which modifies the scale-dependent gravitational potential and induces 
observable signatures in weak-lensing convergence fields. Previous Companions 
developed topological transport statistics, multi-probe likelihoods, neural 
surrogates, hierarchical Bayesian inference, and emulator-accelerated pipelines. 
In this work, we assemble these components into a unified end-to-end RDCC 
inference pipeline. The pipeline begins with simulations, extracts multi-probe 
summary statistics, evaluates neural joint likelihoods, incorporates emulator 
uncertainty, applies hierarchical Bayesian modeling, and produces final 
posteriors for $\alpha$. This Companion provides the complete computational and 
statistical workflow required for applying RDCC to Euclid-like survey data.
\end{abstract}
% ---------------------------------------------------------
\section{Pipeline Overview}

The end-to-end RDCC inference pipeline consists of the following stages:

\begin{enumerate}
    \item Simulation of weak-lensing convergence maps for a grid of $\alpha$.
    \item Extraction of multi-probe summary statistics:
    \begin{itemize}
        \item transport-based persistence-diagram statistics,
        \item Minkowski functionals,
        \item shear power spectra.
    \end{itemize}
    \item Training of neural joint likelihoods (Companion LXXXVI).
    \item Calibration via flow-matching and score diagnostics (Companion LXXXVII).
    \item Emulator acceleration using neural transport surrogates (Companion LXXXVIII).
    \item Hierarchical Bayesian inference with transport-regularized posteriors 
          (Companion LXXXIX).
    \item Posterior sampling and uncertainty quantification.
\end{enumerate}
% ---------------------------------------------------------
\section{Simulation Stage}

We generate weak-lensing convergence maps $\kappa(\theta)$ for a grid of $\alpha$ 
values using:

\begin{itemize}
    \item N-body simulations,
    \item fast approximate solvers,
    \item multi-fidelity simulation flows (Companion LXXXVIII).
\end{itemize}

Each simulation produces:

\begin{itemize}
    \item convergence maps,
    \item persistence diagrams,
    \item Minkowski functional curves,
    \item shear power spectra.
\end{itemize}
% ---------------------------------------------------------
\section{Summary-Statistic Extraction}

We compute:

\begin{itemize}
    \item transport distances $\mathcal{T}_{\ell}$,
    \item Minkowski functionals $V_{i}(\nu)$,
    \item shear power spectra $C_{\ell}^{\gamma\gamma}$.
\end{itemize}

These form the multi-probe summary vector:



\[
    \mathbf{s}
    =
    \left\{
        \mathcal{T}_{\ell},
        V_{0}(\nu),
        V_{1}(\nu),
        V_{2}(\nu),
        C_{\ell}^{\gamma\gamma}
    \right\}.
\]


% ---------------------------------------------------------
\section{Neural Joint Likelihood Evaluation}

We evaluate the neural joint likelihood $\mathcal{L}_{\theta}$ from Companion 
LXXXVI:



\[
    \mathcal{L}_{\theta}(\mathbf{s} \mid \alpha)
    =
    \mathcal{N}\!\left(
        f_{\theta}(\mathbf{s});
        \mu(\alpha),
        \Sigma(\alpha)
    \right)
    \left|
        \det \frac{\partial f_{\theta}}{\partial \mathbf{s}}
    \right|.
\]



Calibration ensures statistical correctness (Companion LXXXVII).
% ---------------------------------------------------------
\section{Hierarchical Bayesian Inference}

We define the hierarchical posterior:



\[
    p(\alpha, \theta_{\mathrm{emu}} \mid \mathbf{s}_{\mathrm{obs}})
    \propto
    \mathcal{L}_{\theta}(\mathbf{s}_{\mathrm{obs}} \mid \alpha, \theta_{\mathrm{emu}})
    p(\theta_{\mathrm{emu}})
    p(\alpha).
\]



Transport regularization (Companion LXXXIX) enforces geometric consistency.
% ---------------------------------------------------------
\section{Posterior Sampling and Final Constraints}

We sample the posterior using:

\begin{itemize}
    \item Hamiltonian Monte Carlo,
    \item NUTS,
    \item variational inference,
    \item neural posterior estimation.
\end{itemize}

The final output is:



\[
    p(\alpha \mid \mathbf{s}_{\mathrm{obs}})
\]



with credible intervals and diagnostic checks.
% ---------------------------------------------------------
\section{Conclusion}

This Companion assembled all previous components into a unified end-to-end RDCC 
inference pipeline. The pipeline integrates simulations, multi-probe summary 
statistics, neural likelihoods, emulator acceleration, hierarchical Bayesian 
modeling, and transport regularization. It provides a complete, scalable, and 
statistically robust workflow for inferring the RDCC relaxation parameter $\alpha$ 
from Euclid-like weak-lensing data.

\bibliographystyle{apsrev4-2}
\nocite{*}
\bibliography{references}


\end{document}
